Compression is normally a storage decision. It buys capacity and bandwidth, and
the bill arrives at read time: before anything can be computed, the data must be
expanded back into the form the application expects. For datasets that have
outgrown memory, it is this materialization step, not the compression itself,
that sets the cost of using the data.

We revisit that trade-off for pangenome graphs, the reference data structure of
modern human genomics, whose interchange format now ships hundreds of gigabytes
per release and is growing faster than the cohorts it encodes. We show that a
grammar-compressed pangenome container can serve as the execution substrate for
the analyses that consume it, with no expansion to text and no whole-graph object
in memory. \sys evaluates six analyses as streaming folds over compressed
traversal columns, paying only for the columns a query touches plus a compact
accumulator. Among them is a reference-relative variant projection, the entry
point to the entire downstream genomics toolchain, which we recast as a two-phase
algorithm that separates a topology-only site-discovery pass from a
compressed-traversal allele fold and so never builds the graph.

The result inverts the usual relationship between compression and access cost:
because the compressed form is what the query reads, a better ratio makes queries
faster rather than slower. Against the established tools \sys is $15$--$613\times$
faster while using $3$--$25\times$ less memory, and it holds peak resident set
below the size of the uncompressed input for four of six analyses. It completes
whole-genome analyses on $358$--$369$\,GB graphs, from a $4.5$\,GB container, that
the graph-resident baselines cannot load at all. We also report where the model
breaks down: one analysis builds an uncompressed join index whose footprint,
unlike the traversal store, does grow with the graph, and it is that accumulator,
not the compressed data, that exhausts memory first.
